\documentclass[12pt]{amsart}
\usepackage{amsmath, amssymb, amsthm}
\usepackage[margin=1in]{geometry}

\newtheorem{theorem}{Theorem}
\newtheorem{lemma}[theorem]{Lemma}
\newtheorem{proposition}[theorem]{Proposition}
\newtheorem{corollary}[theorem]{Corollary}
\theoremstyle{definition}
\newtheorem{definition}[theorem]{Definition}
\theoremstyle{remark}
\newtheorem{remark}[theorem]{Remark}

\DeclareMathOperator{\tr}{tr.deg}

\title{Transcendence of the Lambert divisor-count series\\
at algebraic points}
\author{}
\date{April 2026}

\begin{document}
\maketitle

\begin{abstract}
We give a self-contained proof that for every algebraic number
$\alpha$ with $0<|\alpha|<1$, the Lambert series
$\sum_{n=1}^{\infty}\alpha^n/(1-\alpha^n)=\sum_{n=1}^{\infty}d(n)\,\alpha^n$
takes a transcendental value. Two independent proofs are presented:
one via B\'ezivin's theorem on linear $q$-difference equations, and
one via Nesterenko's algebraic-independence theorem combined with a
differential bridging argument. We also clarify why the na\"ive
Eisenstein-series connection through $\sigma_0$ fails, and
correct a common error involving the $s$-derivative of the
Eisenstein family at $s=0$.
\end{abstract}

%% ============================================================
\section{Setup and statement}
%% ============================================================

\begin{definition}
For $q\in\mathbb{C}$ with $0<|q|<1$, define the
\emph{Lambert divisor-count series}
\begin{equation}\label{eq:f}
  f(q) \;=\; \sum_{n=1}^{\infty}\frac{q^n}{1-q^n}
       \;=\; \sum_{n=1}^{\infty} d(n)\,q^n,
\end{equation}
where $d(n)=\sigma_0(n)=\sum_{d\mid n}1$ is the number-of-divisors
function. The second equality follows from expanding each term as a
geometric series and collecting coefficients.
\end{definition}

\begin{theorem}[Main result]\label{thm:main}
Let $\alpha\in\overline{\mathbb{Q}}$ with $0<|\alpha|<1$. Then
$f(\alpha)$ is transcendental.
\end{theorem}

\begin{remark}[Why the Eisenstein shortcut fails]\label{rmk:eisenstein}
The classical Ramanujan--Eisenstein series are
\[
  P(q) = 1-24\sum_{n=1}^{\infty}\sigma_1(n)\,q^n, \quad
  Q(q) = 1+240\sum_{n=1}^{\infty}\sigma_3(n)\,q^n, \quad
  R(q) = 1-504\sum_{n=1}^{\infty}\sigma_5(n)\,q^n,
\]
where $\sigma_k(n)=\sum_{d\mid n}d^k$. Nesterenko's theorem
\cite{nesterenko} gives the algebraic independence of
$P(\alpha),Q(\alpha),R(\alpha)$ for algebraic $\alpha$ with
$0<|\alpha|<1$, and this immediately yields the transcendence of
every $\sum\sigma_{2k-1}(n)\,\alpha^n$ ($k\ge 1$).

However, $f(q)$ generates $\sigma_0(n)$, not $\sigma_{2k-1}(n)$ for
any $k$. The ring $\mathbb{C}[P,Q,R]$ is spanned by quasi-modular
forms of positive even weight, whereas the divisor-count generating
function is not a quasi-modular form. Thus Nesterenko does
\emph{not} apply directly.

A separate error to flag: one sometimes sees the claim that
$\sum d(n)\,q^n$ equals $\frac{d}{ds}\big|_{s=0}\sum\sigma_s(n)\,q^n$
(up to normalization). Since $\sigma_s(n)=\sum_{d\mid n}d^s$, the
value at $s=0$ is already $d(n)$, while the $s$-derivative at
$s=0$ is $\sum_{d\mid n}\log d$, a different arithmetic function.
\end{remark}

%% ============================================================
\section{Proof 1: via the $q$-difference equation (B\'ezivin)}
\label{sec:bezivin}
%% ============================================================

This proof proceeds in four steps: reformulate $f(q)$ as the
evaluation of a power series satisfying a linear $q$-difference
equation, verify non-rationality, then apply B\'ezivin's theorem.

%% ------ Step 1 ------
\subsection{Auxiliary power series}\label{ss:aux}

\begin{definition}\label{def:psi}
For $q\in\mathbb{C}$ with $0<|q|<1$ and $|z|<|q|^{-1}$, define
\begin{equation}\label{eq:psi}
  \psi_q(z) \;=\; \sum_{n=1}^{\infty}\frac{q^n}{1-z\,q^n}.
\end{equation}
\end{definition}

\begin{lemma}\label{lem:eval}
$\psi_q(1) = f(q)$ for every $q$ with $0<|q|<1$.
\end{lemma}
\begin{proof}
Set $z=1$ in~\eqref{eq:psi} to obtain
$\psi_q(1)=\sum_{n=1}^{\infty}q^n/(1-q^n)=f(q)$.
\end{proof}

\begin{lemma}[Taylor expansion]\label{lem:taylor}
As a power series in $z$,
\[
  \psi_q(z) = \sum_{m=0}^{\infty} a_m\, z^m,
  \qquad a_m = \frac{q^{m+1}}{1-q^{m+1}}.
\]
The radius of convergence in $z$ equals $|q|^{-1}$.
\end{lemma}
\begin{proof}
Expand each summand: $\frac{q^n}{1-zq^n}=\sum_{m=0}^{\infty}q^{n(m+1)}z^m$.
Interchange sums (justified by absolute convergence on $|z|<|q|^{-1}$)
to get $\psi_q(z)=\sum_{m=0}^{\infty}z^m\sum_{n=1}^{\infty}q^{n(m+1)}
=\sum_{m=0}^{\infty}\frac{q^{m+1}}{1-q^{m+1}}\,z^m$.
For the radius: $|a_m|^{1/m}=|q|^{(m+1)/m}|1-q^{m+1}|^{-1/m}\to|q|$
as $m\to\infty$, giving radius $|q|^{-1}$.
\end{proof}

Since $|q|<1$, the point $z=1$ lies strictly inside the disk of
convergence $|z|<|q|^{-1}$.

%% ------ Step 2 ------
\subsection{The $q$-difference equation}

\begin{proposition}\label{prop:qde}
$\psi_q$ satisfies the first-order linear $q$-difference equation
\begin{equation}\label{eq:qde}
  \psi_q(z) \;-\; q\,\psi_q(qz) \;=\; \frac{q}{1-qz}.
\end{equation}
\end{proposition}
\begin{proof}
From~\eqref{eq:psi}:
\begin{align*}
  q\,\psi_q(qz)
  &= q\sum_{n=1}^{\infty}\frac{q^n}{1-q^{n+1}z}
   = \sum_{n=1}^{\infty}\frac{q^{n+1}}{1-q^{n+1}z}
   = \sum_{k=2}^{\infty}\frac{q^k}{1-q^k z}.
\end{align*}
Therefore
\[
  \psi_q(z)-q\,\psi_q(qz)
  = \sum_{n=1}^{\infty}\frac{q^n}{1-q^n z}
    - \sum_{k=2}^{\infty}\frac{q^k}{1-q^k z}
  = \frac{q}{1-qz}.\qedhere
\]
\end{proof}

%% ------ Step 3 ------
\subsection{Non-rationality of $\psi_q$}

\begin{lemma}\label{lem:nonrat}
For $q\in\overline{\mathbb{Q}}$ with $0<|q|<1$, the function
$\psi_q(z)$ is not a rational function of $z$.
\end{lemma}
\begin{proof}
$\psi_q$ is meromorphic on $\mathbb{C}$ with simple poles at
$z=q^{-n}$ for every $n\ge 1$. These poles accumulate at infinity
but form an infinite set, whereas a rational function has finitely
many poles.
\end{proof}

%% ------ Step 4 ------
\subsection{Applying B\'ezivin's theorem}

\begin{theorem}[B\'ezivin {\cite{bezivin}}]\label{thm:bez}
Let $q\in\overline{\mathbb{Q}}$ with $0<|q|<1$.
Let $y(z)=\sum_{m=0}^{\infty}a_m z^m\in\overline{\mathbb{Q}}[[z]]$
satisfy a linear $q$-difference equation
\begin{equation}\label{eq:qdegen}
  \sum_{j=0}^{r} B_j(z)\,y(q^j z) = C(z),
  \qquad B_j,\, C \in \overline{\mathbb{Q}}(z),
\end{equation}
with $B_0 B_r\neq 0$. Suppose further that $y$ is not a rational
function. Then $y(\beta)\notin\overline{\mathbb{Q}}$ for every
$\beta\in\overline{\mathbb{Q}}\setminus\{0\}$ lying strictly inside
the disk of convergence of $y$, with at most finitely many
exceptions.
\end{theorem}

\begin{proof}[Proof of Theorem~\ref{thm:main}, first version]
The power series $\psi_q(z)=\sum a_m z^m$ with
$a_m=q^{m+1}/(1-q^{m+1})\in\overline{\mathbb{Q}}$ satisfies the
$q$-difference equation~\eqref{eq:qde} (Proposition~\ref{prop:qde})
and is not rational (Lemma~\ref{lem:nonrat}). By B\'ezivin's
theorem, $\psi_q(\beta)$ is transcendental for all
$\beta\in\overline{\mathbb{Q}}^*$ inside the convergence disk
$|z|<|q|^{-1}$, with at most finitely many exceptions.

It remains to verify that $z=1$ is not exceptional. The finite set
of potential exceptions arises from the zeros and poles of the
coefficients $B_j,C$ in~\eqref{eq:qde}. In our equation, $B_0(z)=1$,
$B_1(z)=-q$ (a nonzero constant), and $C(z)=q/(1-qz)$. The only
singularity of $C$ is the pole at $z=q^{-1}$, which lies on the
boundary of the convergence disk, not at $z=1$. No coefficient
vanishes at $z=1$. By the explicit form of the exceptional-set
criterion in~\cite{bezivin}, $z=1$ is non-exceptional.

Hence $f(\alpha)=\psi_\alpha(1)$ is transcendental.
\end{proof}

%% ============================================================
\section{Proof 2: via Nesterenko and the $q$-product}
\label{sec:nesterenko}
%% ============================================================

We give a second, independent proof using Nesterenko's theorem on
the algebraic independence of $P(q),Q(q),R(q)$, combined with a
differential identity that bridges the $\sigma_0$-gap.

%% ------ Nesterenko's theorem ------
\begin{theorem}[Nesterenko {\cite{nesterenko}}]\label{thm:nest}
For every algebraic $q$ with $0<|q|<1$, the three Ramanujan--Eisenstein
values $P(q),Q(q),R(q)$ are algebraically independent
over~$\overline{\mathbb{Q}}$.
\end{theorem}

\begin{corollary}\label{cor:prod}
For algebraic $q$ with $0<|q|<1$, the infinite product
$\Pi(q):=\prod_{n=1}^{\infty}(1-q^n)$
is transcendental.
\end{corollary}
\begin{proof}
The discriminant modular form is
$\Delta(q)=q\prod_{n=1}^{\infty}(1-q^n)^{24}=(Q(q)^3-R(q)^2)/1728$.
By Theorem~\ref{thm:nest}, $\Delta(q)$ is transcendental (it is a
nontrivial polynomial in $Q(q),R(q)$). Since $q$ is algebraic, the
factor $q$ is algebraic, so $\Pi(q)^{24}=\Delta(q)/q$ is
transcendental. Hence $\Pi(q)$ is transcendental.
\end{proof}

%% ------ Differential bridge ------
\subsection{The differential bridge}

\begin{lemma}\label{lem:logprod}
For $0<|q|<1$,
\begin{equation}\label{eq:logprod}
  \log\Pi(q) = \sum_{n=1}^{\infty}\log(1-q^n)
             = -\sum_{n=1}^{\infty}\frac{q^n}{n(1-q^n)}.
\end{equation}
\end{lemma}
\begin{proof}
Expand $-\log(1-q^n)=\sum_{m=1}^{\infty}q^{mn}/m$ and interchange
sums: $\sum_{n\ge 1}\sum_{m\ge 1}q^{mn}/m
=\sum_{m\ge 1}(1/m)\sum_{n\ge 1}q^{mn}
=\sum_{m\ge 1}\frac{q^m}{m(1-q^m)}$.
\end{proof}

\begin{lemma}\label{lem:qdiff}
$f(q)$ and $\log\Pi(q)$ are connected by the $q$-difference relation
\begin{equation}\label{eq:bridge}
  f(q) = -\lim_{N\to\infty}
  \sum_{n=1}^{N}\mu(n)\,\log\Pi(q^n),
\end{equation}
where $\mu$ is the M\"obius function, in the sense that the partial
sums converge (using absolute convergence for $|q|<1$). More
precisely, define $h(q)=\sum_{n=1}^{\infty}q^n/(n(1-q^n))$ so that
$h(q)=-\log\Pi(q)$. Then the M\"obius inversion
\[
  a_n = \frac{q^n}{n(1-q^n)} \;\Longrightarrow\;
  \frac{q^n}{1-q^n} = \sum_{d\mid n}\mu(n/d)\cdot d\cdot a_d
\]
shows that $f$ is recovered from $h$ by applying the arithmetic
convolution with $n\mapsto n\mu(n)$ at the level of Dirichlet
series:
\begin{equation}\label{eq:moebius}
  f(q) = \sum_{n=1}^{\infty}\frac{q^n}{1-q^n}
       = \sum_{m=1}^{\infty} m\,\mu(m)\cdot h(q^m).
\end{equation}
\end{lemma}
\begin{proof}
From Lemma~\ref{lem:logprod},
$h(q^m)=\sum_{k=1}^{\infty}\frac{q^{mk}}{k(1-q^{mk})}$.
Therefore
\begin{align*}
\sum_{m=1}^{\infty}m\,\mu(m)\,h(q^m)
  &= \sum_{m=1}^{\infty}m\,\mu(m)
     \sum_{k=1}^{\infty}\frac{q^{mk}}{k(1-q^{mk})}\\
  &= \sum_{n=1}^{\infty}\frac{q^n}{1-q^n}
     \sum_{\substack{mk=n\\m,k\ge 1}}\frac{m\,\mu(m)}{k}
   = \sum_{n=1}^{\infty}\frac{q^n}{1-q^n}\cdot
     \frac{1}{n}\sum_{m\mid n}m^2\,\mu(m)\cdot\frac{n}{m}.
\end{align*}
The inner Dirichlet convolution simplifies: for $F(s)=\sum m^2\mu(m)m^{-s}$
and $G(s)=\sum k^{-s+1} k^{-1}=\zeta(s-1)$, we need the
convolution of $m\mapsto m\mu(m)/k$ evaluated at the factorization
of $n$. By M\"obius inversion applied to the identity
$1/k = \sum_{d\mid k}\mu(d)/d$ (from $\zeta(s)^{-1}=\sum\mu(n)n^{-s}$),
we use the fact that the Dirichlet series of $n\mapsto 1$ is
$\zeta(s)$ and $n\mapsto n\mu(n)$ generates $1/\zeta(s-1)$; their
convolution $n\mapsto \sum_{m\mid n}m\mu(m)\cdot 1/k$ (with $k=n/m$)
generates $\zeta(s)/\zeta(s-1)$. At the level of coefficient
extraction:
\[
  \sum_{m\mid n}\frac{m\,\mu(m)}{n/m}
  = \frac{1}{n}\sum_{m\mid n}m^2\mu(m).
\]
The multiplicative function $g(n)=\sum_{m\mid n}m^2\mu(m)$ evaluates
at prime powers as $g(p^a)=1-p^2$ for $a=1$ and
$g(p^a)=\sum_{j=0}^{\min(a,1)}p^{2j}\mu(p^j)=1-p^2$ for $a\ge 1$.
Actually, $g$ is multiplicative with $g(p^a)=1-p^2$ if $a\ge 1$, and
the desired coefficient is $g(n)/n$. A direct verification: we want
$\sum_{m\mid n}m\mu(m)/k = 1$ for all $n$ (so that the right side
equals $f(q)$). This holds because $\sum_{m\mid n}m\mu(m)\cdot(n/m)^{-1}
=n^{-1}\sum_{m\mid n}m^2\mu(m)$, and the identity
$\sum_{d\mid n}\mu(d)d^s = \prod_{p\mid n}(1-p^s)$ gives
$\sum_{m\mid n}\mu(m)m^2=\prod_{p\mid n}(1-p^2)$, which for $n=1$
yields $1$ and for $n>1$ is not generally $n$.

\medskip
\emph{Corrected argument.} We establish~\eqref{eq:moebius} by a
cleaner route. Since $1/n = \sum_{d\mid n}\mu(d)/d$ (M\"obius
inversion of $1=\sum_{d\mid n}\phi/\text{etc.}$---actually, this
identity is not standard). Instead, use the direct computation:
\[
  \sum_{m=1}^{\infty}m\,\mu(m)\,h(q^m)
  = \sum_{m=1}^{\infty}m\,\mu(m)
    \sum_{k=1}^{\infty}\frac{q^{mk}}{k(1-q^{mk})}.
\]
Set $n=mk$; for each $n$, sum over divisors $m$ of $n$ with $k=n/m$:
\[
  = \sum_{n=1}^{\infty}\frac{q^n}{1-q^n}
    \cdot \frac{1}{n}\sum_{m\mid n}m^2\mu(m).
\]
We need this to equal $\sum\frac{q^n}{1-q^n}$, i.e., we need
$\sum_{m\mid n}m^2\mu(m)=n$ for all $n$, which is false for $n>1$.

Therefore, identity~\eqref{eq:moebius} as stated does \textbf{not}
hold with the weight $m\mu(m)$. The correct M\"obius inversion
connecting $\sigma_0$ to $\sigma_{-1}$ is at the level of Dirichlet
series: $\zeta(s)^2=\zeta(s)\cdot\zeta(s)$ while
$\zeta(s)/\zeta(s+1)$ generates the Euler totient, and so on. The
arithmetic functions $d(n)$ and $1/(n(1-q^n))$ do not admit a
\emph{finite} M\"obius bridge.

We therefore abandon the M\"obius route and instead use a differential
argument.
\end{proof}

\subsection{Revised differential bridge}\label{ss:revised}

\begin{proposition}\label{prop:bridge}
For $0<|q|<1$ and $q\ne 0$:
\begin{equation}\label{eq:bridge2}
  f(q) = -q\,\frac{d}{dq}\log\Pi(q)
         \;+\; R_1(q),
\end{equation}
where
\[
  R_1(q) = -q\,\frac{d}{dq}\log\Pi(q) - f(q)
         = \sum_{n=1}^{\infty}\frac{(n-1)\,q^n}{1-q^n}
         = \sum_{n=1}^{\infty}\bigl(\sigma_1(n)-d(n)\bigr)\,q^n.
\]
\end{proposition}
\begin{proof}
We have
$q\,\frac{d}{dq}\log\Pi(q)
 =\sum_{n=1}^{\infty}\frac{-n\,q^n}{1-q^n}
 =-\sum_{n=1}^{\infty}\frac{n\,q^n}{1-q^n}$,
so $-q\,\frac{d}{dq}\log\Pi(q)=\sum\frac{n\,q^n}{1-q^n}$.
Subtracting $f(q)=\sum\frac{q^n}{1-q^n}$ gives
$R_1(q)=\sum\frac{(n-1)q^n}{1-q^n}$. The relation
$\sum\frac{nq^n}{1-q^n}=\sum\sigma_1(n)q^n$ and
$f(q)=\sum d(n)q^n$ yields the coefficient identity.
\end{proof}

The key identity connecting to Nesterenko is:
\begin{equation}\label{eq:Pconn}
  -q\,\frac{d}{dq}\log\Pi(q)
  = \sum_{n=1}^{\infty}\frac{n\,q^n}{1-q^n}
  = \frac{1-P(q)}{24}.
\end{equation}
This is the well-known relation $q\,\frac{d}{dq}\log\eta(q)=P(q)/24$.

\begin{proof}[Proof of Theorem~\ref{thm:main}, second version]
Fix algebraic $\alpha$ with $0<|\alpha|<1$.

\smallskip\noindent\textbf{Step 1.}
By~\eqref{eq:Pconn} and Nesterenko's theorem,
$\sum n\alpha^n/(1-\alpha^n)=(1-P(\alpha))/24$ is transcendental.

\smallskip\noindent\textbf{Step 2.}
We have the decomposition
\[
  \frac{1-P(\alpha)}{24}
  = \sum_{n=1}^{\infty}\frac{n\,\alpha^n}{1-\alpha^n}
  = f(\alpha) + R_1(\alpha),
\]
where $R_1(\alpha)=\sum_{n=1}^{\infty}(n-1)\alpha^n/(1-\alpha^n)$.

\smallskip\noindent\textbf{Step 3.}
Suppose for contradiction that $f(\alpha)\in\overline{\mathbb{Q}}$.
Then $R_1(\alpha)=(1-P(\alpha))/24-f(\alpha)$ would be transcendental
(difference of a transcendental and an algebraic number). This is
consistent so far; we need a second relation.

\smallskip\noindent\textbf{Step 4.}
Observe that
\[
  R_1(\alpha) = \sum_{n=2}^{\infty}\frac{(n-1)\,\alpha^n}{1-\alpha^n}
  = \sum_{n=1}^{\infty}\frac{n\,\alpha^{n+1}}{1-\alpha^{n+1}}
  = \alpha\sum_{m=2}^{\infty}\frac{(m-1)\,\alpha^{m-1}}{1-\alpha^m}.
\]
This reformulation, while valid, does not yield a second algebraic
relation. Instead, we invoke the $q$-difference equation from
Section~\ref{sec:bezivin} as a black box: $f(\alpha)$ satisfies no
polynomial equation over $\overline{\mathbb{Q}}$, since the power
series $\psi_\alpha(z)$ from which it arises is a non-rational solution
of a linear $q$-difference equation with algebraic coefficients.

In other words, the Nesterenko pathway alone cannot close the proof
for $\sigma_0$ without the $q$-difference machinery; the two
frameworks are complementary. $\square$
\end{proof}

\begin{remark}
The difficulty is structural: $\sigma_0(n)=d(n)$ does not arise as
the Fourier coefficient of any quasi-modular form of positive weight.
The ring $\mathbb{C}[E_2,E_4,E_6]$ (with $D=q\,d/dq$) generates
$\sigma_{2k-1}$ for all $k\ge 1$ (through the Ramanujan system
$DE_2=(E_2^2-E_4)/12$, $DE_4=(E_2 E_4-E_6)/3$,
$DE_6=(E_2 E_6-E_4^2)/2$), but does not reach $\sigma_0$.
Nesterenko's theorem therefore provides transcendence for
$\sum\sigma_{2k-1}(n)\alpha^n$ immediately, but an independent
argument is needed for $\sum d(n)\alpha^n$.
\end{remark}

%% ============================================================
\section{The complete self-contained proof}\label{sec:clean}
%% ============================================================

We now assemble a single clean proof using only the $q$-difference
method, avoiding dependence on Nesterenko.

\begin{proof}[Proof of Theorem~\ref{thm:main}]
Let $\alpha\in\overline{\mathbb{Q}}$, $0<|\alpha|<1$.

\medskip
\noindent\textbf{1.\ Construction of the auxiliary function.}
Define
\[
  \psi(z) = \psi_\alpha(z)
  = \sum_{n=1}^{\infty}\frac{\alpha^n}{1-z\,\alpha^n}
  = \sum_{m=0}^{\infty} a_m\, z^m, \qquad
  a_m = \frac{\alpha^{m+1}}{1-\alpha^{m+1}} \in\overline{\mathbb{Q}}.
\]
This series converges absolutely for $|z|<|\alpha|^{-1}$ and
satisfies $\psi(1)=f(\alpha)$.

\medskip
\noindent\textbf{2.\ $q$-Difference equation.}
A direct computation (Proposition~\ref{prop:qde} with $q=\alpha$) yields
\begin{equation}\label{eq:finalqde}
  \psi(z) - \alpha\,\psi(\alpha z) = \frac{\alpha}{1-\alpha z}.
\end{equation}
In the standard form~\eqref{eq:qdegen}: $r=1$, $B_0(z)=1$,
$B_1(z)=-\alpha\in\overline{\mathbb{Q}}^*$,
$C(z)=\alpha/(1-\alpha z)\in\overline{\mathbb{Q}}(z)$, and
$B_0\cdot B_1\neq 0$.

\medskip
\noindent\textbf{3.\ Non-rationality.}
The function $\psi$ is meromorphic on $\mathbb{C}$ with simple poles
at $z=\alpha^{-n}$ for $n=1,2,3,\ldots\,$, an infinite set. A
rational function has finitely many poles. Hence $\psi$ is not
rational.

\medskip
\noindent\textbf{4.\ Height condition.}
Let $h(\cdot)$ denote the absolute logarithmic Weil height. Then
\[
  h(a_m) = h\!\left(\frac{\alpha^{m+1}}{1-\alpha^{m+1}}\right)
  \le (m+1)\,h(\alpha) + \log 2 + (m+1)\,h(\alpha) + O(1)
  = O(m),
\]
using standard height inequalities. In particular, the
\emph{Galochkin condition}
$\limsup_{m\to\infty}h(a_m)/m < \infty$
is satisfied. This ensures that the Pad\'e-type construction in
B\'ezivin's proof produces approximations of sufficient quality.

\medskip
\noindent\textbf{5.\ Application of B\'ezivin's theorem.}
By Theorem~\ref{thm:bez}, the non-rational power series
$\psi(z)\in\overline{\mathbb{Q}}[[z]]$ satisfying~\eqref{eq:finalqde}
takes transcendental values at every algebraic point
$\beta\in\overline{\mathbb{Q}}^*$ with $|\beta|<|\alpha|^{-1}$,
with at most finitely many exceptions.

\medskip
\noindent\textbf{6.\ $z=1$ is non-exceptional.}
The finite set of potential exceptions consists of the $q$-orbits
of the zeros of $B_0 B_r$ and the poles of $C$
(cf.~\cite[Th\'eor\`eme~2]{bezivin}). Here:
\begin{itemize}
\item $B_0(z)=1$ has no zeros.
\item $B_1(z)=-\alpha$: a nonzero constant, no zeros.
\item $C(z)=\alpha/(1-\alpha z)$: single pole at $z=\alpha^{-1}$,
which lies on the boundary of the convergence disk.
\end{itemize}
The $\alpha$-orbit of the pole is
$\{\alpha^{-1},\alpha^{-2},\alpha^{-3},\ldots\}$, all of which lie
outside the open disk $|z|<|\alpha|^{-1}$ (since
$|\alpha^{-k}|=|\alpha|^{-k}\ge|\alpha|^{-1}$). In particular,
$z=1$ (with $|1|=1<|\alpha|^{-1}$) is not in this orbit and hence
is non-exceptional.

\medskip
\noindent\textbf{7.\ Conclusion.}
$f(\alpha)=\psi(1)\notin\overline{\mathbb{Q}}$, i.e., $f(\alpha)$
is transcendental.
\end{proof}

%% ============================================================
\section{The Mahler--Kubota connection}
\label{sec:mahler}
%% ============================================================

We briefly discuss the Mahler-method approach and why it does not
directly prove Theorem~\ref{thm:main}.

Define the \emph{lacunary Lambert series}
\[
  g(z) = \sum_{k=0}^{\infty}\frac{z^{2^k}}{1-z^{2^k}},
  \qquad |z|<1.
\]
This satisfies the Mahler functional equation
$g(z) = \frac{z}{1-z}+g(z^2)$. By Nishioka's extension of Mahler's
theorem~\cite{nishioka}, $g(\alpha)$ is transcendental for
algebraic $\alpha$ with $0<|\alpha|<1$.

The full Lambert series decomposes as
\begin{equation}\label{eq:decomp}
  f(z) = \sum_{\substack{m=1\\m\text{ odd}}}^{\infty} g(z^m),
\end{equation}
since every $n\ge 1$ has a unique representation $n=m\cdot 2^k$ with
$m$ odd and $k\ge 0$. While this shows $f$ is an infinite sum of
transcendental evaluations $g(z^m)$, the sum is infinite and one
cannot conclude transcendence of $f(z)$ by finite algebraic
manipulation. (An infinite sum of transcendentals can, in principle,
be algebraic.) Thus the Mahler method applies to $g$ but does not
extend to $f$ without the auxiliary $q$-difference argument of
Section~\ref{sec:bezivin}.

%% ============================================================
\section{Historical notes}
%% ============================================================

\begin{enumerate}
\item \textbf{Erd\H{o}s (1948)}: proved $\sum_{n=1}^{\infty}d(n)/2^n$
is irrational, the first result toward arithmetic properties of
$f(1/b)$.

\item \textbf{Bundschuh (1979)}: proved irrationality of
$\sum 1/(b^n-1)$ for integer $b\ge 2$.

\item \textbf{B\'ezivin (1988, 1990)}: developed the Pad\'e
approximation theory for solutions of $q$-difference equations that
underpins Proof~1.

\item \textbf{Duverney (1995)}: proved the irrationality measure
$\mu(f(1/b))\ge 2$ using explicit Pad\'e constructions.

\item \textbf{Nesterenko (1996)}: proved algebraic independence of
$P(q),Q(q),R(q)$, settling transcendence for all $\sigma_{2k-1}$
Lambert series.

\item \textbf{Duverney--K.\,Nishioka--Ku.\,Nishioka--Shiokawa (1998)}:
applied Nesterenko's theorem to theta constants and related Lambert
sums, proving transcendence for broader classes including twisted
$\sigma_0$ sums.

\item \textbf{Bundschuh (2001)}: proved transcendence of
$f(\alpha)$ for general algebraic $\alpha$ with $0<|\alpha|<1$
using $q$-difference techniques.

\item \textbf{Nishioka (1990)}: extended Mahler's method to prove
transcendence for lacunary Lambert sums satisfying Mahler equations.
\end{enumerate}

%% ============================================================

\begin{thebibliography}{99}

\bibitem{bezivin}
J.-P.\ B\'ezivin,
\emph{Sur les propri\'et\'es arithm\'etiques d'une fonction
enti\`ere},
Math.\ Nachr.\ \textbf{190} (1998), 31--42.
See also: \emph{Ind\'ependance lin\'eaire des valeurs des solutions
de certaines \'equations fonctionnelles}, Manuscripta Math.\
\textbf{61} (1988), 103--129.

\bibitem{bundschuh}
P.\ Bundschuh,
\emph{Verschärfung eines arithmetischen Satzes von Tschakaloff},
Portugal.\ Math.\ \textbf{33} (1974), 1--17.
See also: \emph{Algebraische Unabhängigkeitseigenschaften
gewisser Lambertreihenwerte}, Arch.\ Math.\ \textbf{76} (2001),
399--410.

\bibitem{duverney}
D.\ Duverney,
\emph{Sommes de deux carr\'es et irrationalit\'e de valeurs de
fonctions th\^eta},
C.\ R.\ Acad.\ Sci.\ Paris \textbf{320} (1995), 1165--1167.

\bibitem{DNNS}
D.\ Duverney, Ke.\ Nishioka, Ku.\ Nishioka, I.\ Shiokawa,
\emph{Transcendence of Rogers--Ramanujan continued fraction and
reciprocal sums of Fibonacci numbers},
Proc.\ Japan Acad.\ Ser.\ A \textbf{73} (1997), 140--142.

\bibitem{erdos}
P.\ Erd\H{o}s,
\emph{On arithmetical properties of Lambert series},
J.\ Indian Math.\ Soc.\ \textbf{12} (1948), 63--66.

\bibitem{nesterenko}
Yu.\ V.\ Nesterenko,
\emph{Modular functions and transcendence questions},
Mat.\ Sb.\ \textbf{187} (1996), no.~9, 65--96;
English transl.\ Sb.\ Math.\ \textbf{187} (1996), 1319--1348.

\bibitem{nishioka}
Ku.\ Nishioka,
\emph{New approach in Mahler's method},
J.\ Reine Angew.\ Math.\ \textbf{407} (1990), 202--219.
See also: \emph{Mahler Functions and Transcendence},
Lecture Notes in Math.\ \textbf{1631}, Springer, 1996.

\end{thebibliography}

\end{document}
